% Security checklist for professors: handout for the briefing "Cybersecurity in 2026".
% Compile with XeLaTeX. DRAFT: items marked [verify] rest on sources that we
% could not open. source_notes.pdf lists the status of each claim.
\documentclass[10pt]{article}
\usepackage[letterpaper,margin=0.5in]{geometry}
\usepackage{fontspec}
\setmainfont{Libertinus Serif}
\setsansfont{Libertinus Sans}
\setmonofont{Libertinus Mono}[Scale=0.9]
\usepackage{xcolor,enumitem,amssymb,titlesec,hyperref}
\definecolor{brownred}{HTML}{8B1A1A}
\hypersetup{colorlinks=true,allcolors=brownred}
\titleformat{\section}{\large\sffamily\bfseries\color{brownred}}{}{0pt}{}
\titlespacing*{\section}{0pt}{4pt}{1pt}
\setlist[itemize]{label=$\square$,leftmargin=1.6em,itemsep=0.5pt,topsep=2pt,parsep=0pt}
\newcommand{\mins}[1]{\hfill{\footnotesize\color{gray}#1}}
\ifdefined\hideverify\newcommand{\vfy}{\unskip}\else\newcommand{\vfy}{{\footnotesize\color{gray}[verify]}}\fi
\pagestyle{empty}
\setlength{\parindent}{0pt}
\begin{document}\footnotesize

{\LARGE\sffamily\bfseries\color{brownred} Security checklist for professors}\par
\smallskip
Start with level A, and do the steps of each level from top to bottom. The labels (A1, B2) match the slides. We sorted the steps by 2 criteria: how often attackers use the weakness, and how long the step takes. The sort and the times are our estimates. These steps are for your own devices and accounts. For a Brown-owned device or a Brown account, ask OIT first. Menu names change with each software version.\ifdefined\hideverify\else{} A step marked \vfy{} comes from a source that we could not open for this draft.\fi

\section{A. This week (about 2 hours)}
\begin{itemize}
\item \textbf{A1. Updates.} Turn on automatic updates on the laptop, the phone, and the browser. Restart today. \mins{15 min}
\item \textbf{A2. Password manager.} Install one: 1Password, or Bitwarden at no cost (Wirecutter picks \vfy), or Dashlane, or the built-in Apple Passwords or Google Password Manager. The company cannot read your passwords or reset the main password \vfy: write the main password on paper. Replace the passwords for email and the bank with long, unique passwords. \mins{60 min}
\item \textbf{A3. Sign-in methods.} A passkey is a secret that your device stores. It replaces the password and the second step. Google: g.co/passkeys, then ``Create a passkey'' \vfy. Apple: turn on two-factor authentication and iCloud Keychain \vfy. Turn on a second sign-in step (2FA) on every account that offers one. Strength, from most to least: security key, passkey, app code, text-message code. \mins{30 min}
\item \textbf{A4. Phone passcode.} Use at least 6 digits \vfy. iPhone: turn on Stolen Device Protection (Settings, Face ID \& Passcode) \vfy. Android: turn on Theft protection (Settings, Google, All services) \vfy. \mins{10 min}
\item \textbf{A5. Habit 1.} Never approve a sign-in prompt, or enter a sign-in code, that you did not request. Never give a sign-in code to another person.
\item \textbf{A5. Habit 2.} Never follow a web page's instruction to paste text into a terminal or the Windows Run box.
\item \textbf{A5. Habit 3.} When someone asks for money or a password by email, call, or video, call back on a number that you already have. Report suspicious messages to OIT.
\item \textbf{A5. Code word.} Agree on a code word with your family. Ask for it when a caller who sounds like a relative asks for money \vfy. \mins{15 min}
\item \textbf{A6. Phone number lock.} Set a PIN or a number lock with your carrier, in its app or on its website, so that nobody can move your number. Verizon: Number Lock and SIM Protection \vfy. T-Mobile: SIM Protection and Port Out Protection \vfy. AT\&T: Wireless Account Lock \vfy. \mins{5 min}
\end{itemize}

\section{B. This month (about 4 hours)}
\begin{itemize}
\item \textbf{B1. Security key setup.} Print the backup codes first (B2). A security key is a small device that you plug in or tap to sign in. Buy 2 (Wirecutter's pick: Yubico Security Key C NFC, about \$29 each \vfy), add both keys to your Apple and Google accounts in the security settings, and keep the spare key at home. Apple: Settings, your name, Sign-In \& Security, Two-Factor Authentication, Security Keys \vfy. If you lose every key and every trusted device, Apple cannot let you back in \vfy. \mins{45 min}
\item \textbf{B2. Account recovery.} First, print the backup codes, and add a recovery contact at Apple and at Google \vfy. Then remove text-message codes, once 2 other methods work \vfy. Run myaccount.google.com/security-checkup \vfy. \mins{30 min}
\item \textbf{B3. Browser.} Turn on ``Always use secure connections'' (Chrome) or ``HTTPS-Only Mode'' (Firefox) \vfy. Install uBlock Origin (Firefox) or uBlock Origin Lite (Chrome, Edge, Safari), and Privacy Badger (Wirecutter recommendations \vfy). Remove the extensions that you do not use. \mins{15 min}
\item \textbf{B4. Backups.} Keep 1 backup copy of your own files offline \vfy. Restore 1 file as a test. Keep research and student data only on storage that Brown approves. \mins{60 min}
\item \textbf{B4. Disk encryption.} On your own computer: back up first, turn on FileVault (Mac) or BitLocker (Windows), and save the recovery key away from the laptop \vfy. \mins{20 min}
\item \textbf{B5. Credit freeze.} Place a free freeze at each of the 3 credit bureaus (Equifax, Experian, TransUnion) \vfy. Lift it a few days before you apply for a loan or a card. \mins{30 min}
\item \textbf{B5. Old devices.} Replace devices that no longer receive updates, such as an old home router \vfy.
\item \textbf{B6. Old accounts.} Delete the accounts that you no longer use. Find them in your password manager and in old ``welcome'' emails. Download the files that you need first. \mins{30 min}
\end{itemize}

\section{C. When time allows, or at higher risk (personal accounts)}
\begin{itemize}
\item \textbf{C1. Google Advanced Protection Program.} It requires a passkey or a security key, limits other apps' access to Gmail and Drive, and makes recovery slower \vfy. Enroll at landing.google.com/advancedprotection \vfy.
\item \textbf{C1. Apple Advanced Data Protection.} Set a recovery contact or a recovery key first. Then: Settings, your name, iCloud, Advanced Data Protection \vfy. If you lose your devices and that contact or key, Apple cannot recover the data \vfy.
\item \textbf{C1. Apple Lockdown Mode.} For the few people whom sophisticated attackers target. Settings, Privacy \& Security, Lockdown Mode \vfy.
\item \textbf{C2. Email masking.} Hide My Email (Apple, with iCloud+): Settings, your name, iCloud, Hide My Email \vfy. Others: Firefox Relay, DuckDuckGo Email Protection, Fastmail Masked Email, or masked addresses in 1Password \vfy.
\item \textbf{C3. Private Relay.} Apple, with iCloud+, for Safari: Settings, your name, iCloud, Private Relay \vfy. It can fail on campus networks \vfy.
\item \textbf{C3. VPN.} Use the Brown VPN for Brown resources. A commercial VPN hides your traffic from the local network, and the VPN company can then see it \vfy. Wirecutter's pick: [pick].


\item \textbf{C4. Tracking.} iPhone: Settings, Privacy \& Security, Tracking; turn off ``Allow Apps to Request to Track'' \vfy. Android: Settings, Privacy, Ads; choose ``Delete advertising ID'' \vfy. Give apps your location only while you use them. Opt out at the largest data brokers \vfy.
\end{itemize}

\section{Research labs}
\begin{itemize}
\item Keep study data on storage that Brown approves for that kind of data.
\item Give each student a named account, and remove access when a person leaves.
\item Delete data that the study no longer needs, after you check the retention rules of the sponsor, the IRB, and Brown.
\item Tell OIT about instrument computers that run old systems. Do not update or disconnect them yourself.
\item Wait about 7 days before you install a new release of an open-source code package \vfy. Update phones, laptops, and browsers at once.
\item API keys: keep them out of code and prompts, give each project its own key with a spending limit, and limit a server's key to the address of that server.
\item Run AI coding agents in a container. Do not connect an AI agent to Brown email or Drive without approval from OIT.
\end{itemize}

\section{After a suspected break-in}
\begin{itemize}
\item Call OIT first, from another phone. Follow its instructions for the device.
\item Do not erase the device, delete files, or try to repair it.
\item On a different device, change the password of the affected account, and sign out of all other devices.
\item If money moved, call the bank.
\end{itemize}

\end{document}
